\documentclass[10pt,a4paper]{amsart}
\usepackage[margin=19mm]{geometry}
\usepackage{amsmath,amssymb,mathtools}
\usepackage[scr=rsfs,cal=euler]{mathalfa}
\usepackage{xfrac}
\usepackage[unicode,hidelinks,bookmarks=false]{hyperref}
\newcommand{\ie}{i.e.\ }
\newcommand{\cf}{cf.\ }
\newcommand{\resp}{respectively\ }
\newcommand{\Subsection}{\subsection}
\providecommand{\nobreakdash}{\nobreak\hbox{-}}
\setlength{\emergencystretch}{2em}
\multlinegap0pt
\usepackage[shortlabels]{enumitem}
\usepackage{mathtools}
\usepackage{amssymb}
\newtheorem{lemma}{Lemma}
\newtheorem{definition}{Definition}
\newtheorem{proposition}{Proposition}
\DeclareMathOperator{\Hom}{\textup{Hom}}
\DeclareMathOperator{\RPC}{\mathbf{RPC}}
\DeclareMathOperator{\Span}{span}
\def\hom{\mathrm{hom}}
\def\trop{\mathrm{trop}}
\title{Birational Classification of Orbifold Compactified Jacobians}
\author{Jeremy Feusi, Sam Molcho}
\date{Source: arXiv:2605.03835v1 (CC BY 4.0)}
\begin{document}
\maketitle

\noindent\emph{Source and licence.} Birational Classification of Orbifold Compactified Jacobians. Version arXiv:2605.03835v1 (CC BY 4.0), distributed by arXiv under the Creative Commons Attribution 4.0 International licence, http://creativecommons.org/licenses/by/4.0/, as recorded in the arXiv licence metadata for this exact version and shown on the abstract page https://arxiv.org/abs/2605.03835v1. Full licence notice: \texttt{licenses/arxiv-2605.03835-CC-BY-4.0.txt}.\par\smallskip

\noindent\textit{Editorial scope.} Counted unit: the article's proposition prop:cone\_cx\_iff\_fan (source0.tex lines 2337--2346). The article's complete original proof of it is delivered with it (source0.tex lines 2523--2605, one \texttt{proof} environment of 83 lines, which the article places away from the statement). No mathematics is written, repaired or re-derived: the statement and the proof are the article's own lines, in the article's own notation, and no retained line is substituted or re-flowed.  Retained uncounted as the source-local context the proof needs: lines 1080--1084; lines 2293--2312; lines 2350--2354.  Nothing was omitted from any retained span, so the package carries no omission record.  No retained line contains a citation, so the wrapper carries no bibliography and no bibliography entry.  No retained line contains a figure environment, an image-inclusion command or drawing code, so no image is delivered and the package declares no asset dependency.  Editorial formatting only, in addition: an AMS \texttt{amsart} preamble that restates the article's own declarations and macros used by the retained lines, namely the environments \texttt{lemma}, \texttt{definition}, \texttt{proposition} on separate counters, together with the standard packages those lines need.  The retained environments print in this document as Lemma 1, Definition 1, Proposition 1 in order of appearance, whereas the article prints the same items with the numbering of its own full document.  The complete native source of the article as distributed in the arXiv e-print for this exact version is bundled unchanged as \texttt{sources/199820/source0.tex}.
\begin{lemma}
	\label{lem:pos_definite_vanishing}
	Fix a symmetric positive-definite bilinear form $Q$ over $\sigma\in \RPC$.
	Let $m\in M$ and $n\in N_{\sigma}\cap \sigma$. If $n(Q(m,m))=0$ then $n(Q(m,m'))=0$ for all $m'\in M$.
\end{lemma}

\begin{definition}
	\label{def:stacky_fan}
	Let $A:=\Hom(\underline{M},\mathbb{G}_{m,\trop})^{(M)}/Q_{\hom}(\underline{M})$ be a tropical abelian variety
	over an RPC $\sigma$. Assume moreover we are given a second lattice $M'$.

	A stacky fan is an (in general infinite) set $\Delta$ of rational polyhedral cones $\tau\subseteq (N_{\sigma}\times N)^{(M)}\times N'$,
	equipped with finite-index sublattices $\tilde{N}_{\tau}\subseteq ((N_{\sigma}\times N)\times N')\cap \Span(\tau)$
	satisfying the following conditions:
		\begin{enumerate}
			\item Any two elements of $\Delta$ intersect along a common face.
			\item For any $\tau,\tau'\in \Delta$ we have $\tilde{N}_{\tau}\cap \Span(\tau\cap \tau')=\tilde{N}_{\tau'}\cap \Span(\tau\cap \tau')$.
			\item For any $\tau\in \Delta$ and any face $\tau'$ of $\tau$ we have $\tau'\in \Delta$.
			\item For any $\tau\in \Delta$ and $m\in M$, we have $T_m(\tau)\in \Delta$ equipped with the lattices $T_m(\tilde{N}_{\tau})$ (where $T_m$ acts trivially on the $N'$ factor).
			\item For any $\tau\in \Delta$ and $m\in M$, we have that $\tau\cap T_m(\tau)$ is (pointwise) fixed by $T_m$.
			\item There are only finitely many equivalence classes of cones in $\Delta$, where $\tau_1,\tau_2\in \Delta$ are equivalent if and only if $\tau_2=T_m(\tau_1)$ for some $m\in M$.
			\item The cone $\sigma\times\{0\}\times\{0\}\subseteq N_{\sigma}\times N\times N'$ with lattice
				$\tilde{N}_{\sigma\times \{0\}\times\{0\}}:=(N_{\sigma}\times N\times N')\cap \Span(\sigma\times \{0\}\times\{0\})$ is an element of $\Delta$.
		\end{enumerate}
		We say that $\Delta$ is complete if $(N_{\sigma}\times N)^{(M)}\times N'\subseteq \bigcup_{\tau\in \Delta}^{}\tau$.
\end{definition}

\begin{lemma}
	\label{lem:diag_repr}
	The diagonals $T_{N'}^{\trop}\to T_{N'}^{\trop}\times T_{N'}^{\trop}$ and
	$A\to A\times_{\sigma}A$ are representable by finite cone complexes.
\end{lemma}

\begin{proposition}
	\label{prop:cone_cx_iff_fan}
	In the above situation, the following statements hold:
\begin{enumerate}
	\item $\Sigma_{\Delta}$ is representable by a finite cone complex.
	\item If $\tilde{N}_{\tau'}=\Span(\tau')\cap (N_{\sigma}\times N \times N')$ for all $\tau'\in \Delta$ and $\Delta$ is complete, then $\Sigma_{\Delta}\to A\times T_{N'}^{\trop}$ is a subdivision.
	\item The morphism $\Sigma_{\Delta}\to A\times T_{N'}^{\trop}$ is a partial tropical compactification of $A\times T_{N'}^{\trop}\to \sigma$.
	\item Every partial tropical compactification $\Sigma\to A\times T_{N'}^{\trop}$ with $\Sigma$ a finite cone complex is isomorphic to $\Sigma_{\Delta}$ for a unique stacky fan $\Delta$.
\end{enumerate}
\end{proposition}

\begin{proof}[Proof of Proposition \ref{prop:cone_cx_iff_fan}]
	Step 1 (Show that $\Sigma_{\Delta}$ is a finite cone complex): To lighten the notation, we assume that $M'=0$ in this proof. The proof goes through verbatim for the general case.

	Pick a stacky fan $\Delta$. We must show that $\Sigma_{\Delta}$ is representable by a finite cone complex.
	Let $\Sigma$ be the set of equivalence classes of cones in $\Delta$ modulo the equivalence relation in Definition
	\ref{def:stacky_fan}. For each $[\tau]\in \Sigma$ pick a representative $\tau$ and equip $\tau$ with the lattice $\tilde{N}_{\tau}$.
	We define face morphisms $[\tau]\to [\tau']$ to be the datum
	of an RPC $T_m\tau\in [\tau]$ and a morphism
	$$\tau\to T_m\tau\to \tau'$$
	in $\Delta$, where the first
	morphism is the canonical isomorphism and the second is a face inclusion (the element $m$ is not part of the data). This endows
	$\Sigma$ with the structure of a category fibered in setoids over $\RPC$ and it remains to show that $\Sigma$ is a finite cone complex with
	$\Sigma\cong \Sigma_{\Delta}$. The fact that $\Sigma$ has finitely many cones is clear by construction.
	To see that it is a cone complex, first observe that
	if $\tau\in \Sigma$, then for every face $\tau'$ of $\tau$ we have $[\tau']\in \Sigma$ by condition (3) in the definition of a stacky fan. Hence it remains
	to show that for $[\tau_1],[\tau_2]\in \Sigma$, there exists at most one morphism $[\tau_1]\to [\tau_2]$ in $\Sigma$. Assume we are
	given two morphisms $[\tau_1]\to [\tau_2]$ defined by $\tau_1\to T_{m_1}\tau_1\subseteq\tau_2$ and $\tau_1\to T_{m_2}\tau_1\subseteq\tau_2$.
	Then
	$$T_{m_1}\tau_1=T_{m_1-m_2}T_{m_2}\tau_1\subseteq T_{m_1-m_2}\tau_2.$$
	Hence $T_{m_1}\tau_1\subseteq \tau_2\cap T_{m_1-m_2}\tau_2$.
	By condition (5) in the definition of a stacky fan, the element $m_1-m_2\in M$ fixes $\tau_2\cap T_{m_1-m_2}\tau_2$ pointwise,
	so $T_{m_1}\tau_1=T_{m_2}\tau_1$ and it follows that the two morphisms are equal.

	Step 1.1 (Show $\Sigma_{\Delta}\cong \Sigma$): The datum of a morphism $\tau\to \Sigma_{\Delta}$ for $\tau\in \RPC$
	is equivalent to the datum of a morphism $\tau\to N_{\sigma,\mathbb{R}}\times N_{\mathbb{R}}$ such that the image of $\tau$ is contained
	in some $\tau_0\in \Delta$, modulo the action of $M$ on $\Hom(\tau,(N_{\sigma}\times N)^{(M)})$. Translating $\tau\to N_{\sigma}\times N$
	if necessary, we may assume that $\tau_0$ is the chosen representative of $[\tau_0]\in \Sigma$. This defines a morphism $\tau\to \Sigma$,
	independent of the choice of translation of $\tau_0$ by condition (5) in the definition of a stacky fan. Conversely, given $\tau\to \Sigma$ factoring through
	$\tau_0\in \Sigma$, composing with $\tau\to \tau_0\to N_{\sigma,\mathbb{R}}\times N_{\mathbb{R}}$ gives a morphism $\tau\to \Sigma_{\Delta}$. It is clear from
	construction that these two assignments are inverses to one another and we conclude.

	Step 2 (Show $\Sigma_{\Delta}\to A$ is a partial tropical compactification): From Lemma \ref{lem:diag_repr} and the representability of $\Sigma_{\Delta}$, it follows that $\Sigma_{\Delta}\to A$
	is representable by finite cone complexes. We have a section $\sigma\to \Sigma_{\Delta}$ given by the cone $\sigma\times\{0\}\in \Delta$,
	which lies over the zero-section $\sigma\to A$. It is clear from the construction of $\Sigma_{\Delta}$ that $\Sigma_{\Delta}\to A$
	is a monomorphism, hence $\Sigma_{\Delta}\to A$ is a partial tropical compactification of $A\to \sigma$.

	If each $\tilde{N}_{\tau}$ is given by $\Span(\tau)\cap (N_{\sigma}\times N)$ and $\Delta$ is complete, then it follows from the
	above construction that $\Sigma_{\Delta}\to A$ is a subdivision.

	Step 3 (Show that every representable partial tropical compactification is a stacky fan): We show that any partial tropical compactification $\Sigma\to A$ with $\Sigma$ representable by a finite cone complex is isomorphic
	to $\Sigma_{\Delta}$ for a unique stacky fan. For uniqueness, note that if $\Sigma\cong \Sigma_{\Delta}$, then $\Delta$
	can be recovered as the set of cones $\{T_m\tau|\tau\in \Sigma\}$ together with the finite-index sublattices
	$T_m N_{\tau}$. For existence, we let $\Sigma$ be arbitrary. Let $\tau\in \Sigma$. Then $\tau\to \Sigma\to A$ is a monomorphism,
	so choosing a lift, we may view $\tau\subseteq(N_{\sigma}\times N)^{(M)}$. Let $\Delta:=\{T_m\tau|\tau\in \Sigma\}$
	and for each $T_m\tau$ equip it with the sublattice $\tilde{N}_{T_m \tau}=T_m N_{\tau}$. This is well-defined by assumption, since any
	$m$ fixing $\tau$ acts trivially on the span of $\tau$ by assumption. We show that $\Delta$ is a stacky
	fan. First, it is clear that $T_m \tau\subseteq (N_{\sigma}\times N)^{(M)}$ since $\tau$ is. Moreover, if $\tilde{\tau}$ is a
	face of $T_m \tau$, then $\tilde{\tau}=T_m \tau'$ for some face $\tau'$ of $\tau$, so $\Delta$ is closed under taking faces. Next,
	since $\Sigma$ is finite, $\Delta$ has only finitely many cones up to equivalence. Let $\sigma\to \Sigma$ be a morphism lifting
	the zero section, then the composition $\sigma\to \Sigma\to A$ lies in the equivalence class of the map $\sigma\to \sigma\times \{0\}\subseteq (N_{\sigma}\times N)^{(M)}$,
	so $\sigma\times\{0\}\in \Delta$. Next, we show that $T_{m}\tau\cap \tau$ is pointwise fixed by $m$ for any $\tau\in \Delta$.
	For this, let
	$$x\in T_m\tau\cap\tau\cap \tilde{N}_{\tau}\cap \tilde{N}_{T_m\tau}$$
	and consider the morphism $\mathcal{A}^{1}:=(\mathbb{R}_{\ge 0},\mathbb{Z})\to \tau$
	which sends 1 to $x$. The morphisms $\mathcal{A}^{1}\to T_m\tau$ and $\mathcal{A}^{1}\to \tau$ define lifts of this morphism
	to $\Sigma$. Since $\Sigma\to A$ is a monomorphism, it follows that the two lifts agree, so $x$ must be fixed by the action of $T_m$.

	Step 3.1 (Show that cones of $\Delta$ intersect along common faces): We show that $\tau:=T_{m_1}\tau_1\cap T_{m_2}\tau_2$ is a common face of $T_{m_1}\tau_1$
	and $T_{m_2}\tau_2$ and that the sublattices agree on this common face. We prove this by induction on the maximum $d$ of $\dim(\tau_1)$
	and $\dim(\tau_2)$, the case $d=0$ being clear. Equip $\tau$ with the lattice $T_{m_1}N_{\tau_1}\cap T_{m_2}N_{\tau_2}$.
	Then the compositions $\tau\to T_{m_{i}}\tau_{i}\cong \tau_i$ define two lifts of $\tau\to A$ to $\tau\to \Sigma$. Since $\Sigma\to A$
	is a monomorphism, these two lifts must agree. It follows that the image of $\tau$ in $\tau_1$ is glued to the image of $\tau$ in $\tau_2$
	in $\Sigma$. Since the gluing morphisms in $\Sigma$ are face morphisms, the image of $\tau$ is contained in a common face
	$\tau_0\subseteq \tau_i$ of the cones $\tau_i$. It follows that $\tau=T_{m_1}\tau_0\cap T_{m_2}\tau_0$. If $\dim(\tau_0)<d$, we conclude
	by the induction hypothesis. Otherwise, the image of $\tau$ must intersect the interior of the cone of larger dimension which forces
	$\tau_0=\tau_1=\tau_2$. So we must show that $T_{m_1}\tau_0\cap T_{m_2}\tau_0$ is a face of $\tau_0$. Acting by $T_{-m_2}$,
	we reduce to showing this for $T_{m}\tau_0\cap \tau_0$. In this case, by the previous property, we know that $\tau$ is pointwise
	fixed by $m$. Let $(x,n)\in N_{\sigma}\times N$ be a point lying in $T_m\tau_0\cap \tau_0$.
	So $(x,n+Q_{\hom,N}(m)(x))=(x,n)$, i.e., $Q_{\hom,N}(m)(x)=0$. By Lemma \ref{lem:pos_definite_vanishing}, this is equivalent
	to $Q(m,m)(x)=0$. Since $Q(m,m)\in \sigma^{\vee}\subseteq \tau_0^{\vee}$, it follows that the locus $\{(x,n)|Q(m,m)(x)=0\}$
	is a face of $\tau_0$ (resp. $T_{m}\tau_0$). Hence
	$$T_m\tau_0\cap \tau_0=\{(x,n)\in \tau_0|Q(m,m)(x)=0\}$$
	is a common face
	of $\tau_0$ and $T_m\tau_0$. Moreover, we have $Q(m,m)(x)=0$ for all $x$ in the span of $\tau$, so it follows that
	$\Span(\tau)\cap \tilde{N}_{\tau_0}=\Span(\tau)\cap \tilde{N}_{T_m \tau_0}$ finishing the proof that $\Delta$ is a stacky fan.

	Step 3.2 (Show $\Sigma\cong \Sigma_{\Delta}$): By construction, we have a monomorphism $\Sigma\to \Sigma_{\Delta}$
	which sends $\tau\to \Sigma$ to $\tau\to \tau'\subseteq N_{\sigma,\mathbb{R}}\times N_{\mathbb{R}}$ for any cone $\tau'\in \Sigma$ containing the
	image of $\tau$. Conversely, given a morphism $\tau\to \Sigma_{\Delta}$, let $T_m\tau'\in \Delta$ be a cone such
	that $\tau\to N_{\sigma,\mathbb{R}}\times N_{\mathbb{R}}$ factors through $T_m\tau'$, with $\tau'\in \Sigma$. We obtain a morphism $\tau\to \tau'$
	by composing with the canonical isomorphism $T_m\tau'\cong \tau'$. The morphism $\tau\to \tau'\to \Sigma$ is independent of the
	choice of $m$ and $\tau'$ and defines an inverse $\Sigma_{\Delta}\to \Sigma$ to the above morphism.
\end{proof}

\end{document}
